\chapter{Geometric Matching}
\label{ch:geometric_matching}

\section{Comparing Geometric Objects}
\label{sec:matching_intro}

\index{matching problem}Geometric matching is the algorithmic discipline of
deciding how well two geometric objects agree, and of finding the
\emph{transformation} that makes them agree. The objects can be point sets,
polygonal curves, polygonal regions, or surfaces; the transformation is
usually a rigid motion, a similarity, or an affine map; and the agreement is
measured by a distance between the transformed object and its target. The
three ingredients --- objects, transformations, and measures --- determine a
family of problems with very different algorithmic flavors, unified by a
single question: how much computation does it take to align one shape with
another?

\begin{definition}[Matching Instance]\label{def:match_instance}
  A \textbf{matching instance} is a triple $(\mathcal{A}, \mathcal{B},
  \mathcal{T})$, where $\mathcal{A}$ and $\mathcal{B}$ are geometric objects
  (point sets, curves, or polygons) and $\mathcal{T}$ is a family of
  admissible transformations $T$ of the plane or of space. A solution is a
  transformation $T \in \mathcal{T}$ together with a correspondence
  $\pi$ between the features of $T(\mathcal{A})$ and those of
  $\mathcal{B}$, minimizing a prescribed distance $d(T(\mathcal{A}),
  \mathcal{B})$.
\end{definition}

The correspondence $\pi$ is usually a bipartite matching between the
features (points, vertices, segments) of the two objects. When both objects
are point sets of equal size and $\pi$ must be a bijection, the problem is a
\textbf{bipartite matching problem in the geometric plane}; the two classical
objective functions are the \textbf{minimum-weight matching}
(Definition~\ref{def:match_min_weight}), which minimizes the sum of the
matched distances, and the \textbf{bottleneck matching}
(Definition~\ref{def:match_bottleneck}), which minimizes the largest matched
distance. When the correspondence is free to omit parts of either set, the
problem is \textbf{partial matching}, discussed in
Section~\ref{sec:matching_shape}.

A useful taxonomy organizes the field along three axes:
\begin{itemize}
  \item \textbf{Exact versus approximate}: \emph{exact} matching decides
    whether two objects are congruent (Section~\ref{sec:matching_rigid});
    \emph{approximate} matching minimizes a distance such as the Hausdorff
    distance (Section~\ref{sec:matching_hausdorff}) or the Fr\'echet distance
    (Section~\ref{sec:matching_frechet}).
  \item \textbf{Transformation class}: rigid motions (translations,
    rotations, reflections), similarities (rigid motions plus scaling), and
    affine maps. Larger classes make matching easier but the decision less
    discriminative; Sections~\ref{sec:matching_rigid} and
    \ref{sec:matching_hashing} develop this trade-off.
  \item \textbf{Structure of the objects}: point sets admit the
    combinatorial machinery of matchings and of arrangements
    (Sections~\ref{sec:matching_shape} and \ref{sec:matching_registration});
    curves require order-respecting measures such as the Fr\'echet distance
    and its discrete variant (Sections~\ref{sec:matching_frechet} and
    \ref{sec:matching_discrete_frechet}).
\end{itemize}

\index{matching problem}The chapter proceeds from measures to algorithms.
Sections~\ref{sec:matching_hausdorff}--\ref{sec:matching_discrete_frechet}
define the two principal similarity measures for shapes and curves.
Sections~\ref{sec:matching_shape} and \ref{sec:matching_registration} develop
matching over point sets, including partial matching, outlier handling, the
assignment problem, and the Hungarian algorithm
\cite{kuhn1955,cormen2009}. Sections~\ref{sec:matching_rigid} and
\ref{sec:matching_hashing} treat exact congruence and the model-based
geometric hashing scheme. The iterative closest point algorithm of
Section~\ref{sec:matching_icp} and the applications of
Section~\ref{sec:matching_vision} close the chapter. Throughout, the
robustness questions raised in Section~\ref{sec:math_norms_metrics} of the
mathematical foundations and the nearest-neighbor machinery of
Chapter~\ref{ch:proximity} recur as the computational primitives.

\index{bipartite matching}\index{minimum-weight matching}\index{bottleneck matching}

\section{Hausdorff Distance}
\label{sec:matching_hausdorff}

\index{Hausdorff distance}The \textbf{Hausdorff distance} is the most
important measure of proximity between two sets, and it is the canonical
criterion for approximate point-set matching: it measures the largest
distance a point of either set has to travel to meet the other set.

\begin{definition}[Directed Hausdorff Distance]\label{def:match_directed_hausdorff}
  Let $P$ and $Q$ be two nonempty compact sets of $\mathbb{R}^d$. The
  \textbf{directed Hausdorff distance} from $P$ to $Q$ is
  \[
    h(P, Q) = \max_{p \in P} \min_{q \in Q} \|p - q\| .
  \]
\end{definition}

\begin{definition}[Undirected Hausdorff Distance]\label{def:match_hausdorff}
  The \textbf{Hausdorff distance} between $P$ and $Q$ is the maximum of the
  two directed distances:
  \[
    H(P, Q) = \max\{ h(P,Q),\; h(Q,P) \} .
  \]
\end{definition}

\begin{theorem}[Hausdorff Distance is a Metric]\label{thm:match_hausdorff_metric}
  Restricted to nonempty compact sets, $H$ is a metric: it is nonnegative,
  symmetric, vanishes exactly on equal sets, and satisfies the triangle
  inequality $H(P, R) \le H(P,Q) + H(Q,R)$.
\end{theorem}

\begin{proof}
  Symmetry and nonnegativity are immediate. If $H(P,Q) = 0$ then every point
  of $P$ is at distance $0$ from $Q$, so $P \subseteq Q$; symmetrically
  $Q \subseteq P$, hence $P = Q$. For the triangle inequality, fix $p \in P$
  and let $q \in Q$ attain $\min_{q} \|p - q\|$; then for any $r \in R$,
  \[
    \|p - r\| \le \|p - q\| + \|q - r\| ,
  \]
  and taking the minimum over $r \in R$ and then the maximum over $p \in P$
  gives $h(P,R) \le h(P,Q) + h(Q,R)$; the symmetric inequality is identical
  and $H$ satisfies the triangle inequality.
\end{proof}

For point sets the directed distance is computable from the nearest-neighbor
structure of the target set.

\begin{theorem}[Computing the Hausdorff Distance]\label{thm:match_hausdorff_complexity}
  Let $P$ and $Q$ be point sets with $|P| = n$ and $|Q| = m$. The directed
  distance $h(P, Q)$ can be computed in $O((n + m) \log m)$ time by
  constructing the Voronoi diagram of $Q$ and performing a nearest-neighbor
  query for each $p \in P$ (Section~\ref{sec:proximity_queries} of
  Chapter~\ref{ch:proximity}); the undirected distance $H(P,Q)$ therefore
  takes $O((n + m) \log(n + m))$ time.
\end{theorem}

\begin{proof}
  The Voronoi diagram of $Q$ can be constructed in $O(m \log m)$ time and
  supports nearest-neighbor queries in $O(\log m)$ time each
  (Chapter~\ref{ch:proximity} and \cite{preparata1985,deberg2008}). Querying
  each point of $P$ costs $O(n \log m)$ in total, giving the bound. For the
  undirected distance the roles of $P$ and $Q$ are swapped and the maximum of
  the two directed values is taken.
\end{proof}

\index{Hausdorff distance}\index{min-Hausdorff matching}When a transformation
is allowed, the \textbf{minimum Hausdorff matching} problem asks for the
transformation minimizing the distance after transformation.

\begin{definition}[Minimum Hausdorff Matching]\label{def:match_min_hausdorff}
  Let $P$, $Q$ be point sets and $\mathcal{T}$ a family of transformations.
  The \textbf{minimum Hausdorff distance} is
  \[
    \min_{T \in \mathcal{T}} H(T(P), Q) .
  \]
  A transformation attaining the minimum is a \textbf{min-Hausdorff
  matching} of $P$ into $Q$.
\end{definition}

For translations, the objective function is minimized on an arrangement. Fix
$T$ to be a translation by the vector $t$. For each $p \in P$ the directed
distance to $Q$ is
\[
  \phi_p(t) = \min_{q \in Q} \|p + t - q\| ,
\]
which is the distance from the translated point $p + t$ to its nearest
neighbor in $Q$. The maximum over $p$ of the functions $\phi_p$ is an
\emph{upper envelope of Voronoi surfaces} \cite{alt1999}, and the minimum
Hausdorff distance under translations is the minimum value of this envelope.
The envelope is piecewise defined over the arrangement of the bisectors of
the pairs $(p, q)$ \cite{alt1999}:

\begin{theorem}[Min-Hausdorff under Translations]\label{thm:match_min_hausdorff_translation}
  For point sets $P$ and $Q$ with $|P| = n$ and $|Q| = m$, the minimum of
  $H(t + P, Q)$ over translations $t$ is attained at a vertex of the
  arrangement induced by the distance functions $\phi_p$, and the optimum can
  be found in $O((nm)^2 \log(nm))$ time in the worst case.
\end{theorem}

\begin{proof}[Proof sketch]
  Each function $\phi_p$ is the pointwise minimum of the $m$ functions
  $\|p + t - q\|$, each of which is a translate of the origin-centered
  distance function; the maximum over $p$ is an upper envelope whose
  breakpoints lie on bisectors of the form $\|p + t - q\| = \|p' + t - q'\|$.
  Over any cell of the arrangement of these $O(nm)$ bisectors the envelope is
  a fixed linear distance function, so the minimum over each cell is at a
  cell vertex; a sweep over the arrangement evaluates all vertices. See
  \cite{alt1999} for the exact construction and the refined bounds for
  convex polygons.
\end{proof}

\begin{remark}[Measures for matching]
  The Hausdorff distance does \emph{not} respect the ordering of points on
  curves, which is why curve matching uses the Fr\'echet distance of
  Section~\ref{sec:matching_frechet}. It is also sensitive to outliers,
  since a single stray point of $P$ dominates $h(P,Q)$; the partial matching
  measures of Section~\ref{sec:matching_shape} repair this.
\end{remark}

\index{Hausdorff distance}\index{min-Hausdorff matching}\index{Voronoi surface}

\section{Fr\'echet Distance}
\label{sec:matching_frechet}

\index{Fr\'echet distance}The Hausdorff distance compares curves as
\emph{sets} of points and ignores their parameterization. The
\textbf{Fr\'echet distance} measures instead how far two curves are when a
point moves along one while a leash, held at the other, stays taut: the
distance is the minimal leash length over all monotone ways of walking both
curves from start to finish \cite{frechet1906}. Formally, for polygonal
curves $P = (p_1, \dots, p_n)$ and $Q = (q_1, \dots, q_m)$, the Fr\'echet
distance is
\[
  d_F(P,Q) = \min_{\alpha,\beta \in [0,1]^{[0,1]}} \max_{t} \| P(\alpha(t)) - Q(\beta(t)) \| ,
\]
where the minimum is over all continuous monotone reparameterizations
$\alpha, \beta$ of the two curves.

The Fr\'echet distance is the correct similarity measure for trajectories,
GPS tracks, and any ordered data, precisely because it enforces that the
correspondence between the two curves is order-preserving: two curves that
run in opposite directions are far apart even if their point sets coincide.
The full algorithmic treatment --- the dog-leash analogy, the free-space
diagram, the decision procedure, and the $O(nm \log nm)$ algorithms for
polygonal curves \cite{alt1995} --- is developed in
Section~\ref{sec:frechet_distance} of Chapter~\ref{ch:distance_paths}, to
which we refer. This chapter mentions the measure only to place it in the
matching landscape: for curves it plays the role that the Hausdorff distance
plays for point sets, and it is the natural objective for the ordered point
matching of Section~\ref{sec:matching_shape}.

\index{Fr\'echet distance}\index{free-space diagram}

\section{Discrete Fr\'echet Distance}
\label{sec:matching_discrete_frechet}

\index{discrete Fr\'echet distance}For polygonal curves given by their vertex
sequences, the continuous Fr\'echet distance can be approximated by its
\textbf{discrete} variant, in which the walker and the dog may jump only
between consecutive vertices. The discrete Fr\'echet distance is defined by
monotone couplings of the two vertex sequences \cite{eiter1994}; see
Definition~\ref{def:discrete_frechet} of Chapter~\ref{ch:distance_paths} for
the formal statement.

\begin{definition}[Discrete Fr\'echet Distance, Grid Form]\label{def:match_discrete_frechet_grid}
  Let $P = (p_1, \dots, p_n)$ and $Q = (q_1, \dots, q_m)$ be polygonal
  curves. A \textbf{coupling} is a monotone path in the $n \times m$ grid
  from $(1,1)$ to $(n,m)$ moving only right, down, or diagonally down-right.
  The \textbf{cost} of a coupling $\{(i_k, j_k)\}_{k}$ is
  $\max_k \|p_{i_k} - q_{j_k}\|$, and the \textbf{discrete Fr\'echet
  distance} is the minimum cost over all couplings.
\end{definition}

The grid formulation yields a dynamic program. Let $D(i,j)$ be the discrete
Fr\'echet distance between the prefixes $(p_1, \dots, p_i)$ and $(q_1,
\dots, q_j)$. The recurrence is
\[
  D(i,j) = \max\bigl( \|p_i - q_j\| ,\; \min\{ D(i-1,j),\; D(i,j-1),\;
  D(i-1,j-1) \} \bigr) ,
\]
with $D(1,1) = \|p_1 - q_1\|$, which computes the distance in $O(nm)$ time
and space by filling the grid one diagonal at a time; the coupling achieving
the optimum is recovered by tracing the min-edges backward. The decision
version --- whether the distance is at most $\varepsilon$ --- is the
reachability question in the $\varepsilon$-version of the grid, in which a
cell is \emph{free} when $\|p_i - q_j\| \le \varepsilon$, exactly the
free-space diagram of Section~\ref{sec:frechet_distance}. This is the point
where the discrete and continuous theories meet: the discrete distance
bounds the continuous one from above, and the free-space diagram turns both
into grid reachability problems \cite{alt1995,eiter1994}.

\index{discrete Fr\'echet distance}\index{coupling}

\section{Shape Matching}
\label{sec:matching_shape}

\index{shape matching}\index{partial matching}Shape matching is the general
problem of deciding, for two shapes $\mathcal{A}$ and $\mathcal{B}$, how
similar they are and finding the aligning transformation. Point-set shape
matching subsumes three decisions: which transformation family to admit,
which distance to minimize, and whether all points of one set must be
matched. This section treats the distance and the partial-matching
dimensions; the transformation dimension is developed in
Sections~\ref{sec:matching_rigid}--\ref{sec:matching_hashing}.

\index{bottleneck matching}The two standard objectives for point-set matching
are the minimum-weight matching and the bottleneck matching of the underlying
bipartite distance graph.

\begin{definition}[Minimum-Weight Matching]\label{def:match_min_weight}
  Let $P$ and $Q$ be point sets with $|P| = |Q| = n$. The
  \textbf{minimum-weight matching} is a bijection $\pi \colon P \to Q$
  minimizing
  \[
    \sum_{p \in P} \|p - \pi(p)\| .
  \]
\end{definition}

\begin{definition}[Bottleneck Matching]\label{def:match_bottleneck}
  The \textbf{bottleneck matching} of two point sets of equal size is a
  bijection $\pi$ minimizing the maximum matched distance
  \[
    \max_{p \in P} \|p - \pi(p)\| .
  \]
  Equivalently, it is the smallest radius $r$ for which the bipartite graph
  of pairs $(p,q)$ with $\|p - q\| \le r$ admits a perfect matching.
\end{definition}

The bottleneck version admits a clean decision procedure: sort all
$O(n^2)$ pairwise distances, and find the smallest threshold $r$ for which
the graph $G_r$ has a perfect matching, by a binary search over the sorted
distances combined with a bipartite matching algorithm
\cite{preparata1985,cormen2009}. The minimum-weight version is the
assignment problem of Section~\ref{sec:matching_registration}, solved by the
Hungarian algorithm in $O(n^3)$ time \cite{kuhn1955,cormen2009}. When the
points of $P$ and $Q$ are \emph{ordered} (polygonal curves, sequences of
landmarks), the bijection is constrained to be order preserving and the
problem becomes the sequence matching solved by dynamic programming, with
the Fr\'echet distance of Sections~\ref{sec:matching_frechet} and
\ref{sec:matching_discrete_frechet} as the natural cost
\cite{alt1995,eiter1994}.

\index{outliers}\index{partial Hausdorff distance}Real data rarely consists
of two cleanly aligned point sets: there are \textbf{outliers} (points that
belong to one set but have no counterpart in the other) and
\textbf{clutter}. Partial matching handles this by allowing the matching to
skip points.

\begin{definition}[Partial Matching and Partial Hausdorff]\label{def:match_partial_matching}
  Let $P$ and $Q$ be point sets. A \textbf{partial matching} matches a
  subset $P' \subseteq P$ to $Q$; points of $P \setminus P'$ are
  \textbf{outliers} and pay no cost. The \textbf{partial directed Hausdorff
  distance} $h_k(P,Q)$ is the $k$-th smallest of the distances
  $\min_{q \in Q} \|p - q\|$ over $p \in P$: all but the worst $|P| - k$
  outliers are ignored.
\end{definition}

For the partial distance $h_k(P,Q)$, the outlier sensitivity of the Hausdorff
distance disappears: a set with $|P| - k$ outliers has the same partial
distance as if the outliers were removed. The partial variants of the
minimum-weight and bottleneck matchings, in which a prescribed number of
pairs are dropped, are solved by the same assignment machinery with dummy
rows and columns (Section~\ref{sec:matching_registration}).

\index{matching via the arrangement}Matching also connects to the theory of
arrangements. Fix a translation $t$ and consider the pairs
$(p,q) \in P \times Q$; the translation $t$ maps $p$ within distance $r$ of
$q$ exactly when $t$ lies in the disk $D(p,q,r) = \{ t : \|p + t - q\| \le
r \}$. A perfect matching of radius $r$ exists exactly when the translation
graph defined by these disks has a perfect matching. Varying $r$ and $t$
sweeps an arrangement-like structure --- the same bisector arrangement that
underlies Theorem~\ref{thm:match_min_hausdorff_translation} --- so the
min-Hausdorff, bottleneck, and min-weight problems for translations are all
reducible to searches over the cells of one arrangement
\cite{alt1999}. This is the sense in which ``matching via the arrangement''
unifies the algorithms of this section with the level and envelope machinery
of Chapter~\ref{ch:envelopes} (Section~\ref{sec:arrangement_levels}).

\index{shape matching}\index{partial matching}\index{outliers}\index{bottleneck matching}

\section{Point-Set Registration}
\label{sec:matching_registration}

\index{registration}Point-set registration is the workhorse of the
applications: given two point clouds representing the same scene --- two
scans of a surface, two satellite images, two MRI volumes --- find the
transformation that aligns them, using the pointwise correspondence as an
intermediate structure.

\begin{definition}[Registration]\label{def:match_registration}
  Let $P$ and $Q$ be point sets and $\mathcal{T}$ a family of
  transformations. A \textbf{registration} of $P$ into $Q$ is a
  transformation $T \in \mathcal{T}$ together with a partial matching
  $\pi$ of $T(P)$ into $Q$ minimizing a cost such as
  \[
    \sum_{p} \|T(p) - \pi(p)\|^2
    \qquad\text{or}\qquad
    \max_{p} \|T(p) - \pi(p)\| ,
  \]
  over all admissible $T$ and all matchings $\pi$.
\end{definition}

The optimization over $T$ and $\pi$ jointly is what makes registration hard:
each is easy once the other is fixed. Given the correspondence $\pi$, the
best rigid transformation $T$ is computed by the least-squares alignment of
Section~\ref{sec:matching_rigid}; given $T$, the best correspondence is the
nearest-neighbor matching of each transformed point, i.e., a
\emph{minimum-weight matching} in the bipartite distance graph. The two
optimizations are intertwined, and every practical algorithm --- the
iterative closest point method of Section~\ref{sec:matching_icp} and its
many variants --- alternates between the two, an instance of
alternating-optimization heuristics.

When the correspondence is a bijection and the cost is additive, the
correspondence problem is the classical \textbf{assignment problem}.

\begin{definition}[Assignment Problem]\label{def:match_assignment_problem}
  Given a square matrix $C = (c_{ij})$ of costs $c_{ij} = \|p_i - q_j\|^p$,
  the \textbf{assignment problem} asks for a permutation $\pi$ of
  $\{1,\dots,n\}$ minimizing $\sum_i c_{i\pi(i)}$.
\end{definition}

\begin{theorem}[Hungarian Algorithm]\label{thm:match_hungarian}
  The assignment problem is solvable in $O(n^3)$ time and $O(n^2)$ space by
  the \textbf{Hungarian algorithm} (Kuhn--Munkres), which maintains a set of
  dual variables and finds a perfect matching among the zero-cost edges of
  the transformed cost matrix \cite{kuhn1955,cormen2009}.
\end{theorem}

\begin{algorithm}[htbp]
\caption{Hungarian Algorithm for the Assignment Problem}
\label{alg:match_hungarian}
\begin{algorithmic}[1]
\Require an $n \times n$ cost matrix $C = (c_{ij})$.
\Ensure a permutation $\pi$ minimizing $\sum_i c_{i\pi(i)}$.
\Function{Hungarian}{$C$}
  \State subtract the row minimum from each row, then the column minimum from each column
  \State $\pi \gets \emptyset$ \Comment{partial assignment of zero-cost edges}
  \While{$\pi$ is not a perfect assignment}
    \State mark the rows and columns reachable by alternating zero-cost paths from the unassigned rows
    \State $\Delta \gets \min\{ c_{ij} : i \text{ unmarked}, j \text{ marked} \}$
    \State subtract $\Delta$ from the unmarked rows; add $\Delta$ to the marked columns
    \State augment $\pi$ along a zero-cost augmenting path, if one exists
  \EndWhile
  \State \Return $\pi$
\EndFunction
\end{algorithmic}
\end{algorithm}

In geometric settings the costs $c_{ij} = \|p_i - q_j\|$ are Euclidean
distances, which gives the problem exploitable structure: the closest-pair
and nearest-neighbor machinery of Chapter~\ref{ch:proximity} can prune the
dense cost matrix, and greedy nearest-neighbor matching provides the
excellent initialization that makes the alternating algorithms of
Section~\ref{sec:matching_icp} converge fast in practice. The bottleneck
variant of the assignment problem (minimize $\max_i c_{i\pi(i)}$, cf.\
Definition~\ref{def:match_bottleneck}) is solved by binary search over
thresholds with a bipartite matching feasibility test, in time
$O(n^{2.5} \log n)$ for dense graphs \cite{cormen2009}. Both the min-weight
and the bottleneck assignment problems are used as the correspondence
step of registration pipelines.

\index{assignment problem}\index{Hungarian algorithm}\index{registration}

\section{Rigid Transformations}
\label{sec:matching_rigid}

\index{rigid transformation}Exact matching asks whether two point sets are
\emph{congruent}: whether one can be moved onto the other by an isometry.
The isometries of the plane are the translations, rotations, and reflections
(together with their compositions), and they form the \textbf{Euclidean
group} $E(2)$; in space the analogous group is $E(3)$.

\begin{definition}[Rigid Transformation and Isometry]\label{def:match_isometry}
  A map $T \colon \mathbb{R}^d \to \mathbb{R}^d$ is \textbf{rigid} if it
  preserves all distances: $\|T(x) - T(y)\| = \|x - y\|$ for all $x, y$. The
  set of rigid transformations is the group of \textbf{isometries}. Each
  isometry is a composition of an orthogonal map and a translation; it is
  \textbf{orientation-preserving} (a rotation composed with a translation)
  or \textbf{orientation-reversing} (composed with a reflection).
\end{definition}

Two point sets are congruent if there is an isometry mapping one onto the
other.

\begin{definition}[Congruence]\label{def:match_congruence}
  Point sets $P$ and $Q$ are \textbf{congruent} if there is an isometry $T$
  with $T(P) = Q$ as sets. If $T$ is restricted to translations, the
  appropriate notion is that of \emph{translational congruence}; adding
  rotations gives \emph{rotational congruence}.
\end{definition}

\index{congruence}\index{symmetry}Deciding congruence is fast because
isometries preserve distances. The classical approach sorts the points by
distance to the centroid and matches invariants:

\begin{theorem}[Congruence Testing of Point Sets]\label{thm:match_congruence_test}
  Two point sets $P$ and $Q$ of size $n$ can be tested for congruence in
  $O(n \log n)$ time: after translating both centroids to the origin, the
  multiset of squared distances $\{\|p\|^2 : p \in P\}$ is an invariant, and
  the angular order of the points on their convex hulls
  (Section~\ref{sec:convex_hull_peeling}) resolves the remaining rotations
  and reflections.
\end{theorem}

\begin{proof}[Proof sketch]
  Any isometry fixing the centroid and preserving distances must map the
  points of $P$ to the points of $Q$ at equal distances from the origin, so
  the sorted distance lists must agree. The points at the maximum distance
  (the outer hull vertices) are permuted by the isometry, and the cyclic
  order of the neighbors on the hull fixes the rotation angle; comparing the
  resulting cyclic signatures decides congruence. The sorting dominates, so
  the bound is $O(n \log n)$ \cite{atallah1985symmetry,highnam1986}.
\end{proof}

Symmetries are the congruence tests where $P = Q$: a \emph{symmetry} of a
point set is an isometry mapping the set onto itself. Detecting all
symmetries of a planar point set, and the rotational symmetries of a
polygon, are solved by the same invariant machinery, with optimal
$O(n \log n)$ algorithms \cite{atallah1985symmetry,wolter1985,highnam1986}.

The exact congruences serve two purposes in this chapter. First, they are
the decision problem underlying all approximate matching: an approximate
matching with distance $0$ is a congruence, so the complexity of exact
matching is a lower benchmark for the approximate algorithms. Second, the
least-squares alignment used inside the iterative closest point algorithm
(Section~\ref{sec:matching_icp}) is itself an exercise in rigid
transformations: given a correspondence, the best rigid map is computed
from the singular value decomposition of the cross-covariance of the
matched points, a closed-form computation that is the linear-algebra core of
every registration pipeline.

\index{rigid transformation}\index{isometry}\index{congruence}\index{symmetry}

\section{Iterative Closest Point}
\label{sec:matching_icp}

\subsection{1. Overview}

\textbf{Iterative Closest Point} (ICP) is the dominant algorithm for
registering two point clouds under a rigid transformation. It alternates two
cheap steps --- compute the closest-point correspondences under the current
transform, then compute the rigid transform that best aligns the
correspondences --- and iterates to convergence. First proposed by Besl and
McKay \cite{besl1992} and developed for range-image modeling by Chen and
Medioni \cite{chen1992}, ICP is the standard engine behind surface
reconstruction from range scans
(Section~\ref{sec:shape_reconstruction_point_clouds}) and object
registration in computer vision.

\subsection{2. Definitions}

\begin{definition}[ICP Correspondence]\label{def:match_icp}
  Let $P$ be a \textbf{source} point set and $Q$ a \textbf{target} point
  set. Given a transform $T$, the \textbf{ICP correspondence} matches each
  $p \in P$ to its nearest neighbor in $T(Q)$ (or in $Q$, depending on the
  convention), i.e., to a point $\pi(p)$ minimizing $\|T(p) - \pi(p)\|$.
\end{definition}

\begin{definition}[Alignment Error]\label{def:match_icp_error}
  The \textbf{least-squares alignment error} of a transform $T$ with respect
  to a correspondence $\pi$ is
  \[
    E(T, \pi) = \sum_{p \in P} \|T(p) - \pi(p)\|^2 .
  \]
\end{definition}

\subsection{3. Theory}

The convergence of ICP rests on the fact that each half of the alternation
can only decrease the error.

\begin{theorem}[Monotone Convergence of ICP]\label{thm:match_icp_convergence}
  Let $E_k$ be the alignment error after iteration $k$ of the closest-point
  step followed by the least-squares rigid alignment step. Then the sequence
  $E_k$ is nonincreasing, and the algorithm converges to a local minimum of
  the registration objective \cite{besl1992}.
\end{theorem}

\begin{proof}[Proof sketch]
  In the closest-point step the correspondence is recomputed so that each
  matched target point is at least as close as before, so the error with the
  \emph{old} transform cannot increase; in the alignment step the rigid
  transform minimizing the error for the \emph{fixed} correspondence cannot
  increase it either. The error sequence is thus nonincreasing and bounded
  below, hence convergent. The limit is a local, not global, optimum, which
  is why ICP initialization and its variants matter (Section~5).
\end{proof}

\subsection{4. Pseudocode}

\begin{algorithm}[htbp]
\caption{Iterative Closest Point}
\label{alg:match_icp}
\begin{algorithmic}[1]
\Require source point set $P$, target point set $Q$, initial transform $T_0$.
\Ensure a rigid transform $T$ approximately minimizing the registration error.
\Function{ICP}{$P$, $Q$, $T_0$}
  \State $T \gets T_0$
  \Repeat
    \State build a kd-tree or Voronoi query structure for $Q$
    \For{each $p \in P$}
      \State $\pi(p) \gets$ nearest neighbor of $T(p)$ in $Q$ \Comment{closest-point step}
    \EndFor
    \State $T \gets$ least-squares rigid alignment of $\{(p, \pi(p))\}$
           \Comment{via the singular value decomposition}
  \Until{the decrease in alignment error is below a tolerance}
  \State \Return $T$
\EndFunction
\end{algorithmic}
\end{algorithm}

\subsection{5. Parameter Selection}

\begin{itemize}
\item \textbf{Initialization}: ICP converges only locally
  (Theorem~\ref{thm:match_icp_convergence}), so $T_0$ must already be close;
  coarse registration by matching invariants (centroids, principal axes,
  geometric hashing of Section~\ref{sec:matching_hashing}) is standard.
\item \textbf{Distance threshold}: pairs farther than a threshold are
  rejected as outliers before the alignment step; a threshold of a few
  multiples of the target resolution is typical.
\item \textbf{Point selection}: uniform downsampling, or sampling the most
  salient points (high curvature), speeds the algorithm up with little loss
  of accuracy.
\item \textbf{Query structure}: the nearest-neighbor queries dominate, so
  the kd-tree of Section~\ref{sec:kd_tree} (Chapter~\ref{ch:spatial_structures})
  or the Voronoi structure of Chapter~\ref{ch:proximity} is essential.
\item \textbf{Point-to-plane variants}: replacing the point-to-point cost by
  the distance to the tangent plane at $\pi(p)$ (Chen--Medioni
  \cite{chen1992}) yields faster convergence on smooth surfaces.
\end{itemize}

\subsection{6. Complexity}

\begin{itemize}
\item \textbf{Per iteration}: $O(|P| \log |Q|)$ for the nearest-neighbor
  queries with a kd-tree (Section~\ref{sec:rangesearch_kd_trees}), plus
  $O(|P|)$ for the closed-form alignment via the singular value
  decomposition.
\item \textbf{Iterations}: the number of iterations until convergence is
  data dependent (often tens), so the total cost is $O(k\,|P| \log |Q|)$
  with $k$ the iteration count.
\end{itemize}

\subsection{7. Usages}

\begin{itemize}
\item \textbf{Range-image registration}: aligning consecutive scans of a
  surface into a single model
  \cite{besl1992,chen1992,deberg2008} (Section~\ref{sec:shape_reconstruction_point_clouds}).
\item \textbf{Computer vision}: object tracking, surgical registration, and
  the rigid part of many simultaneous localization and mapping pipelines
  (Section~\ref{sec:matching_vision}).
\item \textbf{Global optimization}: as the local refinement engine inside
  global registration schemes that add loops to the pairwise alignments.
\end{itemize}

\subsection{8. Testing Methods and Edge Cases}

\begin{itemize}
\item \textbf{Synthetic alignment}: generate $Q$ as a known rigid image of
  $P$ with noise, run ICP, and verify the recovered transform.
\item \textbf{Local minima}: test with adversarial initializations far from
  the optimum and verify the algorithm stops at a local, not global, optimum.
\item \textbf{Outliers and partial overlap}: ICP should be tested when only
  a fraction of the points overlap, verifying the threshold rejection.
\item \textbf{Symmetry}: objects with rotational symmetries can converge to
  an alternative valid alignment; verify the error, not the parameters.
\item \textbf{Degenerate clouds}: collinear or coplanar point sets make the
  alignment step ill-conditioned; tests should exercise the SVD-based
  fallback.
\end{itemize}

\subsection{9. References}

\begin{itemize}
\item Besl, P.~J., McKay, N.~D. \cite{besl1992} introduced ICP and proved
  its monotone convergence.
\item Chen, Y., Medioni, G. \cite{chen1992} developed the point-to-plane
  variant for range images.
\item Alt, H., Guibas, L.~J. \cite{alt1999} survey registration within the
  broader matching literature.
\end{itemize}

\index{iterative closest point}\index{ICP}\index{least-squares alignment}

\section{Geometric Hashing}
\label{sec:matching_hashing}

\subsection{1. Overview}

\textbf{Geometric hashing} \cite{wolfson1997} is a model-based object
recognition scheme that finds an object in a cluttered scene without
explicitly searching the transformation space. The idea is to preprocess a
library of models so that \emph{invariant} properties of small point
configurations are stored in a hash table, and then to recognize a scene by
voting: every minimal ``basis'' configuration of scene points casts votes for
the models whose invariants match. The scheme is famous for being
transformation-invariant (rigid, similarity, or affine, depending on the
invariant used) and robust to clutter and partial occlusion, at the price of
a statistical trade-off between the hash bin size and the false-positive
rate.

\subsection{2. Definitions}

\begin{definition}[Geometric Hashing Scheme]\label{def:match_geometric_hashing}
  A \textbf{geometric hashing scheme} consists of
  \begin{itemize}
    \item a family of \textbf{bases}: ordered tuples of points defining a
      coordinate frame (two points for the planar rigid/similarity case,
      three for the affine case);
    \item a \textbf{normalizing map} $F_B$ that maps the basis $B$ onto a
      canonical frame (e.g., the unit segment or unit triangle);
    \item a \textbf{hash function} that quantizes the normalized
      coordinates of the remaining points.
  \end{itemize}
\end{definition}

\subsection{3. Theory}

The correctness of geometric hashing rests on invariance: if the scene is a
rigid (or similarity, or affine) image of a model, then normalizing by the
corresponding basis maps the rest of the model to the same canonical
coordinates regardless of the transformation, so the same hash entries are
reproduced. A scene basis votes for exactly those model bases whose
canonical coordinates coincide under the quantization.

\begin{theorem}[Voting Correctness]\label{thm:match_hashing_voting}
  If the scene contains a subset congruent to a model under the admissible
  transformations, then the basis realizing that congruence accumulates the
  full set of votes of the model, one per matched point, and any threshold
  below the model size separates it from spurious matches with high
  probability when the bins are chosen randomly.
\end{theorem}

\begin{proof}[Proof sketch]
  For the true basis, normalization is exact up to quantization, so every
  scene point that is the image of a model point hashes to the model's bin;
  the number of votes equals the number of matched points. A random model
  point hashes into a given bin with probability proportional to the bin
  area, so the vote count of a wrong model is a binomial random variable
  whose mean is small when the bins are fine. Choosing the bin size and
  threshold to separate the means gives the claim; see \cite{wolfson1997,
  alt1999}.
\end{proof}

\subsection{4. Pseudocode}

\begin{algorithm}[htbp]
\caption{Geometric Hashing (Recognition)}
\label{alg:match_geometric_hashing}
\begin{algorithmic}[1]
\Require a library of models $M_1, \dots, M_K$; a scene point set $S$; a
       basis size $b$ (2 for rigid/similarity, 3 for affine).
\Ensure the models present in $S$ and the aligning transformations.
\Function{GeometricHashingPreprocess}{$M_1, \dots, M_K$}
  \State $\text{HT} \gets$ empty hash table
  \For{each model $M_i$}
    \For{each ordered basis $B$ of $M_i$}
      \State $F \gets$ normalizing map of $B$
      \For{each remaining model point $p$}
        \State $\text{HT}[ \operatorname{quant}(F(p)) ] \gets
               \text{HT}[ \cdot ] \cup \{ (i, B) \}$
      \EndFor
    \EndFor
  \EndFor
  \State \Return $\text{HT}$
\EndFunction
\Function{GeometricHashingRecognize}{$S$, $\text{HT}$, $\tau$}
  \State $\text{votes} \gets$ counter over model--basis pairs
  \For{each ordered basis $B'$ of $S$}
    \State $F' \gets$ normalizing map of $B'$
    \For{each remaining scene point $p$}
      \State increment $\text{votes}[(i, B)]$ for all $(i, B)$ in
             $\text{HT}[ \operatorname{quant}(F'(p)) ]$
    \EndFor
    \State report the pairs $(i, B)$ with $\text{votes} \ge \tau$ and the
           aligning transform $F'^{-1} \circ F$
  \EndFor
\EndFunction
\end{algorithmic}
\end{algorithm}

\subsection{5. Parameter Selection}

\begin{itemize}
\item \textbf{Basis size}: two points for rigid and similarity matching,
  three for affine matching \cite{wolfson1997}; the basis is the 
  ``minimal index'' that determines the transformation class.
\item \textbf{Bin size}: the quantization of the canonical coordinates sets
  the tolerance of the matching; too coarse a bin merges distinct
  configurations (false positives), too fine a bin rejects valid matches.
\item \textbf{Voting threshold $\tau$}: must exceed the expected vote count
  of a wrong model-basis pair; $\tau$ is set from the target false-positive
  rate by the binomial analysis of Theorem~\ref{thm:match_hashing_voting}.
\end{itemize}

\subsection{6. Complexity}

\begin{itemize}
\item \textbf{Preprocessing}: for a model with $n$ points and basis size
  $b$, there are $O(n^b)$ bases and $O(n)$ remaining points each, so
  $O(n^{b+1})$ expected hash entries per model.
\item \textbf{Recognition}: for a scene with $m$ points, $O(m^{b+1})$
  hash lookups in the worst case; with balanced hash tables each lookup is
  expected $O(1)$.
\item \textbf{Space}: the hash table stores $O(\sum_i n_i^{b+1})$ entries.
\end{itemize}

\subsection{7. Usages}

\begin{itemize}
\item \textbf{Model-based object recognition}: recognizing known objects in
  cluttered images despite occlusion and noise \cite{wolfson1997}.
\item \textbf{Coarse registration}: geometric hashing provides the
  transformation-invariant initialization that ICP
  (Section~\ref{sec:matching_icp}) refines, combining global search with
  local precision.
\item \textbf{Biometry}: fingerprint and face recognition pipelines use
  hashed invariant point configurations.
\end{itemize}

\subsection{8. Testing Methods and Edge Cases}

\begin{itemize}
\item \textbf{Occlusion}: hide a fraction of the model points in the scene
  and verify that the vote count tracks the visible fraction.
\item \textbf{Clutter}: add random points to the scene and measure the
  false-positive rate against the binomial prediction of
  Theorem~\ref{thm:match_hashing_voting}.
\item \textbf{Transformation classes}: test rigid, similarity, and affine
  scenes to verify the correct basis size and invariant.
\item \textbf{Degenerate bases}: collinear or coincident basis points give
  degenerate frames; they must be rejected during preprocessing.
\item \textbf{Bin sensitivity}: sweep the bin size and verify the
  precision--recall trade-off is monotone as predicted.
\end{itemize}

\subsection{9. References}

\begin{itemize}
\item Wolfson, H.~J., Rigoutsos, I. \cite{wolfson1997} give the standard
  overview of geometric hashing, its invariants, and its applications.
\item Alt, H., Guibas, L.~J. \cite{alt1999} place geometric hashing in the
  general theory of transformation-invariant matching.
\item The invariant and symmetry machinery of Section~\ref{sec:matching_rigid}
  provides the underlying congruence theory.
\end{itemize}

\index{geometric hashing}\index{voting}\index{invariant}

\section{Applications in Computer Vision}
\label{sec:matching_vision}

\index{computer vision}Geometric matching is the computational core of shape
alignment and registration across vision, medical imaging, and geographic
information systems. The tools of this chapter --- Hausdorff and Fr\'echet
distances, min-weight and bottleneck matchings, the Hungarian algorithm, ICP,
and geometric hashing --- each serve a distinct application profile.

\index{shape alignment}\textbf{Shape alignment} in vision takes a
two-dimensional template and seeks its occurrence in an image. The template
and the image feature points are matched by a min-weight matching or by a
min-Hausdorff matching under rigid and similarity transformations
(Sections~\ref{sec:matching_shape} and \ref{sec:matching_hausdorff}); the
Hausdorff criterion is favored when the image is cluttered, and the
order-respecting Fr\'echet distance
(Section~\ref{sec:matching_frechet}) is used for contour and trajectory
matching. When the transformation class is large or the object can be
occluded, the voting scheme of geometric hashing
(Section~\ref{sec:matching_hashing}) replaces the exhaustive search.

\index{registration}Registration pipelines \cite{besl1992,chen1992,
alt1999} align point clouds acquired at different times, from different
sensors, or from different viewpoints: range scans of an object are merged
into a single model
(Section~\ref{sec:shape_reconstruction_point_clouds} of
Chapter~\ref{ch:shape_reconstruction}), successive frames are stitched into
maps, and medical volumes from different modalities are brought into a
common coordinate frame. The workhorse is ICP
(Section~\ref{sec:matching_icp}) initialized by invariant matching or
geometric hashing, with the assignment problem of
Section~\ref{sec:matching_registration} providing the correspondence step in
feature-based pipelines. In all of these applications the geometric setting
matters: the nearest-neighbor queries inside every iteration are answered by
the kd-trees and Voronoi structures of Chapters~\ref{ch:spatial_structures}
and \ref{ch:proximity}, and robustness to outliers is provided by the
partial matching measures of Section~\ref{sec:matching_shape}.

\index{GIS}\textbf{Geographic information systems} (Chapter~\ref{ch:applications_gis})
use the same machinery to align maps and imagery: registering a satellite
image to a vector map, conflation of road networks from different sources,
and change detection between surveys. Here the transformations are usually
rigid or similarity transforms of image patches, the features are detected
road crossings, building corners, or coastlines, and the matchings are the
min-Hausdorff and bottleneck matchings of
Sections~\ref{sec:matching_hausdorff}--\ref{sec:matching_shape}; the
Fr\'echet distance (Section~\ref{sec:matching_frechet} and
Chapter~\ref{ch:distance_paths}) is the measure of choice for aligning the
ordered sequences of points along roads and rivers.

Across all applications two recurring principles emerge. First,
\emph{initialization and refinement} split the problem: a robust global
method (matching invariants, geometric hashing) produces a coarse alignment
that a local method (ICP) refines. Second, the \emph{measure determines the
algorithm}: outlier-prone data demands the partial and bottleneck measures,
ordered data demands the Fr\'echet family, and clean, fully overlapping data
can use the fastest least-squares and min-weight machinery.

\index{computer vision}\index{shape alignment}\index{registration}\index{GIS}

\section{Exercises}
\label{sec:matching_exercises}

\begin{enumerate}
  \item Prove Theorem~\ref{thm:match_hausdorff_metric} in full detail,
    including the triangle inequality for the undirected Hausdorff distance.
  \item Implement the computation of the Hausdorff distance between two
    point sets (Theorem~\ref{thm:match_hausdorff_complexity}) using a kd-tree
    (Section~\ref{sec:kd_tree}) and verify the running time on random inputs.
  \item Show that the partial Hausdorff distance $h_k(P,Q)$
    (Definition~\ref{def:match_partial_matching}) equals the directed
    Hausdorff distance of the subset of $P$ obtained by deleting the $|P|-k$
    points with the largest nearest-neighbor distances, and use this to give
    an $O((n + m) \log m)$ algorithm for it.
  \item Implement the Hungarian algorithm
    (Algorithm~\ref{alg:match_hungarian}) for the assignment problem and
    verify it against brute force for $n \le 8$ with random cost matrices.
    Then apply it to the Euclidean bottleneck matching
    (Definition~\ref{def:match_bottleneck}) via binary search over the
    pairwise distances.
  \item Prove that the discrete Fr\'echet distance grid recurrence
    (Section~\ref{sec:matching_discrete_frechet}) is correct, and show how
    the free-space diagram decides the $\varepsilon$-reachability question.
  \item Simulate the iterative closest point algorithm
    (Algorithm~\ref{alg:match_icp}) on two congruent copies of a small point
    set under a known rotation and translation, and verify monotone
    convergence of the alignment error (Theorem~\ref{thm:match_icp_convergence}).
  \item For a point set with a rotational symmetry, show that the min-Hausdorff
    distance of Section~\ref{sec:matching_hausdorff} over rotations is
    attained at an angle that is a multiple of the symmetry's angle, and
    test the claim with geometric hashing
    (Section~\ref{sec:matching_hashing}).
  \item Compare, experimentally, the Hausdorff, bottleneck, and min-weight
    matching criteria on a synthetic pair of point sets with $10\%$ outliers
    (Definition~\ref{def:match_partial_matching}): which criterion is most
    robust, and by how much does the running time of the Hungarian algorithm
    (Algorithm~\ref{alg:match_hungarian}) grow with $n$?
\end{enumerate}

\index{exercises}
